\subsection{Mechanistic Interpretability in Transformers}
Elhage et al.\ \cite{Elhage2021} formalised the residual stream framework: the residual stream acts as a central communication channel, with attention heads routing information between tokens and MLP blocks processing within-token representations. Understanding this residual communication is crucial for identifying which layers and components drive specific model behaviours. Building on this, Conmy et al.\ \cite{Conmy2023} introduced ACDC (Automated Circuit Discovery), which identifies the minimal computational subgraph responsible for a target behaviour via iterative activation patching---establishing a rigorous methodology for causal analysis that is, in principle, architecture-agnostic.

\subsection{Sparse Dictionary Learning and Monosemanticity}
A significant obstacle in mechanistic interpretability is \textit{polysemanticity}---individual neurons activating in response to multiple, unrelated features---attributed to \textit{superposition}, where models encode more features than they have dimensions \cite{Elhage2021}. SAEs applied to MLP activations or residual stream vectors learn a larger dictionary of sparse, monosemantic features that influence model predictions in predictable ways \cite{Cunningham2024, Bricken2023}. These contributions target language models exclusively: tokens are words, features correspond to linguistic concepts, and verifying interpretability amounts to reading the activating tokens.

\subsection{Internal Representations of Vision Transformers}
Raghu et al.\ \cite{Raghu2021} showed that ViTs aggregate global information early via self-attention and propagate features strongly through residual connections---properties with no direct analogue in language models that complicate circuit-level analysis. DINO \cite{Caron2021} trains the same ViT architecture through self-distillation, without any label supervision. The resulting representations develop qualitatively different emergent properties---most notably, attention maps that spontaneously segment salient objects---indicating that the training regime shapes internal feature structure substantially. A SAE applied to DINO's MLP activations would therefore likely yield a different vocabulary of sparse features than the texture- and concept-oriented dictionary we observe in the supervised setting. Our codebase already supports DINOv2 as an alternative backbone, but a systematic comparison across training regimes remains a direction for future work.

\subsection{Concept Discovery in Vision-Language Models}
Gandelsman et al.\ \cite{Gandelsman2024} decompose CLIP image representations as sparse sums of text-direction embeddings, yielding human-interpretable visual decompositions. Haque et al.\ \cite{Haque2026} apply SAEs to medical VLMs, anchoring discovered concepts to clinical ontologies and evaluating them with an LLM as external judge. Both approaches rest on pre-existing text-image alignment: the visual concepts are interpretable precisely because they live in a space where text directions are meaningful. In a ViT trained without language supervision, no such space exists---and it is exactly this challenge that this project addresses.
